
\section{\BookRecordIndex{T0333}唯一延拓}\label{sec:2-3}
\zhindex{weiyiyantuo}{唯一延拓}

本节建立 \BookRecordIndex{T0003}$J$-全纯曲线的一个\BookRecordIndex{T0333}唯一延拓定理。在可积情形，这是每个单复变量全纯函数都可作
幂级数展开这一事实的明显推论。在殆复情形，证明困难得多。我们的处理遵循
Floer--Hofer--Salamon \srcite{119} 的工作；另见 Sikorav \srcite{379}。

由于问题是局部的，可以取 $u$ 的定义域为 $\mathbb C$ 中的圆盘
\[
  B_\varepsilon:=\{z\in\mathbb C\mid |z|<\varepsilon\},
\]
并取 $M=\mathbb R^{2n}$。于是\BookRecordIndex{T0088}殆复结构 $J$ 是一个定义在 $\mathbb R^{2n}$ 上、取值于
$\mathbb R^{2n}$ 线性自同构空间 $\GL(2n,\mathbb R)$ 的函数。为供后文引用，在关于 $J$
尽可能弱的光滑性假设下证明\BookRecordIndex{T0333}唯一延拓定理比较方便。因此考虑满足方程
\begin{equation}\label{eq:2-3-1}
  \partial_su+J(u)\partial_tu=0
\end{equation}
的函数 $u:B_\varepsilon\to\mathbb R^{2n}$，其中
$J:\mathbb R^{2n}\to\GL(2n,\mathbb R)$ 是满足 $J^2=-\mathbf1$ 的 $C^1$ 函数。

\begin{remark}\label{rem:2-3-1}
以 $\Delta$ 表示标准\BookRecordIndex{T0220}拉普拉斯算子
\[
  \Delta=(\partial_s)^2+(\partial_t)^2.
\]
利用
$\partial_t(J^2)=(\partial_tJ)J+J(\partial_tJ)=0$，可证明方程
\eqref{eq:2-3-1} 的每个 $W^{1,p}$ 解 $u$（其中 $p>2$）也是二阶拟线性方程
\begin{equation}\label{eq:2-3-2}
  \Delta u=\bigl(\partial_tJ(u)\bigr)\partial_su
  -\bigl(\partial_sJ(u)\bigr)\partial_tu
\end{equation}
的\BookRecordIndex{T0295}弱解。\BookRecordIndex{T0330}椭圆自举论证随即表明：若 $J$ 属于 $C^\ell$，其中 $\ell\ge1$，则每条属于
$W^{1,p}$、$p>2$ 的 \BookRecordIndex{T0003}$J$-全纯曲线 $u$ 必属于 $W^{\ell+1,p}$（见定理 B.4.1 和
注 B.4.3）。特别地，$u$ 属于 $C^1$，所以第 2.2 节的结果适用于这些曲线。
\end{remark}

若可积函数 $w:B_\varepsilon\to\mathbb C^n$ 满足
\[
  \int_{|z|\le r}|w(z)|=O(r^k)
  \qquad\text{对每个 }k\in\mathbb N,
\]
就称 $w$ 在 $z=0$ 处无穷阶消失。若 $w$ 光滑，这表示 $w$ 的无穷阶 jet 在零点消失。
不过，这一条件的积分形式对每个可积函数都有意义。

\begin{theorem}[\BookRecordIndex{T0333}唯一延拓]\label{thm:2-3-2}
设 $u,v\in C^1(B_\varepsilon,\mathbb R^{2n})$ 对某个 $C^1$ \BookRecordIndex{T0088}殆复结构
$J:\mathbb R^{2n}\to\GL(2n,\mathbb R)$ 满足 \eqref{eq:2-3-1}，并且 $u-v$ 在零点
无穷阶消失。则 $u\equiv v$。
\end{theorem}

\begin{corollary}\label{cor:2-3-3}
设 $\Sigma$ 是连通黎曼曲面，$J$ 是光滑流形 $M$ 上的 $C^1$ \BookRecordIndex{T0088}殆复结构。若两条
\BookRecordIndex{T0003}$J$-全纯曲线 $u,v:\Sigma\to M$ 在一点 $z\in\Sigma$ 无穷阶相同，则 $u\equiv v$。
\end{corollary}

\begin{proof}
由定理 2.3.2，$\Sigma$ 中所有使 $u,v$ 无穷阶相同的点组成的集合既开又闭。
\end{proof}

\paragraph{Aronszajn 定理。}
定理 2.3.2 中的函数 $w:=v-u$ 对每个 $p<\infty$ 都属于 $W^{2,p}$（见注 2.3.1），
并在零点无穷阶消失。特别地，$w$ 连续可微。由于 $J$ 及其导数有界，从
\eqref{eq:2-3-2} 可知 $w$ 满足形如下式的微分不等式：
\begin{equation}\label{eq:2-3-3}
  |\Delta w(z)|\le c\bigl(|w|+|\partial_sw|+|\partial_tw|\bigr)
\end{equation}
对每个 $z=s+it\in B_\varepsilon$ 成立。因此，定理 2.3.2 可由下面引用的 Aronszajn
\BookRecordIndex{T0333}唯一延拓定理\BookRecordIndex{T0326}推出。

\begin{theorem}[Aronszajn]\label{thm:2-3-4}
设 $\Omega\subset\mathbb C$ 是连通开集。假设函数
$w\in W^{2,2}_{\mathrm{loc}}(\Omega,\mathbb R^m)$ 在 $\Omega$ 中几乎处处满足逐点估计
\eqref{eq:2-3-3}，并且 $w$ 在某一点 $z_0\in\Omega$ 无穷阶消失。则 $w\equiv0$。
\end{theorem}

这里 $W^{2,2}_{\mathrm{loc}}(\Omega,\mathbb R^m)$ 是这样的映射组成的\BookRecordIndex{T0319}索伯列夫空间：
其二阶导数在 $\Omega$ 的每个预紧开子集上属于 $L^2$。（不熟悉\BookRecordIndex{T0319}索伯列夫空间的读者可以
假定 $w$ 二次连续可微；当 $J$ 属于 $C^2$ 时，这一点自动满足。）Aronszajn 定理可看成
解析函数\BookRecordIndex{T0333}唯一延拓定理的推广。Aronszajn 在 \srcite{18} 中、Hartman 和 Wintner 在
\srcite{169} 中证明了它；比较第 E.4 节。下面给出定理 2.3.2 的另一种证明，它不依赖
Aronszajn 定理。本节余下内容初读时可以略过。

\paragraph{\BookRecordIndex{T0013}Carleman 相似原理。}
考虑线性一阶椭圆系统
\begin{equation}\label{eq:2-3-4}
  \partial_su(z)+J(z)\partial_tu(z)+C(z)u(z)=0.
\end{equation}
这里 $z=s+it$，$\partial_s$ 表示 $\partial/\partial s$，$u$ 是函数
$B_\varepsilon\to\mathbb R^{2n}$，而 $J,C$ 是 $B_\varepsilon$ 上取值于实线性变换空间
$\End(\mathbb R^{2n})=\mathbb R^{2n\times2n}$ 的函数。还假设对每个
$z\in B_\varepsilon$ 都有 $J(z)^2=-\mathbf1$。由于可以取
$J(z)=J'(u(z))$，其中 $J'$ 是 $\mathbb R^{2n}$ 上的\BookRecordIndex{T0088}殆复结构，方程
\eqref{eq:2-3-1} 的任意解都是某个适当函数 $J(z)$ 所对应的方程
\eqref{eq:2-3-4} 的解。不过要注意，即使原来的\BookRecordIndex{T0088}殆复结构 $J'$ 在 $\mathbb R^{2n}$ 上
属于 $C^\ell$，复合映射 $z\mapsto J(z)$ 的光滑性也不会高于 $u$。所以只假设函数
$z\mapsto J(z)$ 属于 $W^{1,p}$。从 \eqref{eq:2-3-1} 转到这里时并不出现含 $C$ 的
线性项，但应用中会需要它。

当我们希望研究方程 \eqref{eq:2-3-1} 的某个给定解 $u$ 的分析性质时，用相应线性方程
\eqref{eq:2-3-4} 代替非线性方程 \eqref{eq:2-3-1} 是合适的。不过，它显然不能用于研究
\eqref{eq:2-3-1} 整个解空间的性质，因为这个线性方程依赖于解 $u$。

对一般非光滑 $J$，理解 \BookRecordIndex{T0003}$J$-全纯曲线\BookRecordIndex{T0408}正则性性质的最有用工具是 \BookRecordIndex{T0013}Carleman 相似原理。
它断言，乘以适当的矩阵值函数 $z\mapsto\Phi(z)$，可把 \eqref{eq:2-3-4} 的任意解 $u$
变换为一个全纯函数 $\sigma$。

\begin{theorem}\label{thm:2-3-5}
固定常数 $p>2$。设
$C\in L^p(B_\varepsilon,\End(\mathbb R^{2n}))$，且
$J\in W^{1,p}(B_\varepsilon,\GL(2n,\mathbb R))$ 满足 $J^2=-\mathbf1$。假设
$u\in W^{1,p}(B_\varepsilon,\mathbb R^{2n})$ 是 \eqref{eq:2-3-4} 的一个解，并且
$u(0)=0$。则存在 $\delta\in(0,\varepsilon)$、映射
\[
  \Phi\in W^{1,p}\bigl(B_\delta,\Hom_{\mathbb R}(\mathbb C^n,\mathbb R^{2n})\bigr)
\]
和全纯函数 $\sigma:B_\delta\to\mathbb C^n$，使 $\Phi(z)$ 可逆，并且对每个
$z\in B_\delta$ 都有
\[
  u(z)=\Phi(z)\sigma(z),
  \qquad \sigma(0)=0,
  \qquad \Phi(z)^{-1}J(z)\Phi(z)=i.
\]
\end{theorem}

从几何上说，可以把映射 $z\mapsto\Phi^{-1}(z)$ 看成丛
$(B_\delta\times\mathbb R^{2n},J)$ 的一个平凡化：它把纤维
$z\times\mathbb R^{2n}$ 上随 $z$ 变化的\BookRecordIndex{T0088}殆复结构族 $J(z)$ 化为常值族 $i$，并把截面
$\widetilde u:z\mapsto(z,u(z))$ 变成全纯截面
$\widetilde\sigma:z\mapsto(z,\sigma(z))$。于是有交换图
\[
\begin{array}{ccc}
  (B_\delta\times\mathbb R^{2n},J)
  &\xrightarrow{\ \Phi^{-1}\ }&(B_\delta\times\mathbb C^n,i)\\[2pt]
  \widetilde u\ \uparrow&&\uparrow\ \widetilde\sigma\\[-1pt]
  B_\delta&\xrightarrow{\quad=\quad}&B_\delta.
\end{array}
\]

\begin{proof}[定理 2.3.2 的证明]
对 $z=s+it\in B_\varepsilon$，定义
$w:B_\varepsilon\to\mathbb R^{2n}$ 为 $w(z):=u(z)-v(z)$。由假设，$w$ 连续可微并在
零点无穷阶消失。此外，$w$ 满足
\begin{align*}
  &\partial_sw(z)+J(u(z))\partial_tw(z)\\
  &\quad=\bigl(J(v(z))-J(u(z))\bigr)\partial_tv(z)\\
  &\quad=\left(\int_0^1\frac{d}{d\tau}
    J\bigl(u(z)+\tau(v(z)-u(z))\bigr)\,d\tau\right)\partial_tv(z)\\
  &\quad=-\left(\int_0^1dJ\bigl(u(z)-\tau w(z)\bigr)w(z)\,d\tau\right)
    \partial_tv(z)\\
  &\quad=-C(z)w(z),
\end{align*}
其中实线性映射 $C(z)\in\End(\mathbb R^{2n})$ 定义为
\[
  C(z)\xi:=
  \left(\int_0^1dJ\bigl(u(z)-\tau w(z)\bigr)\xi\,d\tau\right)\partial_tv(z),
\]
这里 $z\in B_\varepsilon$ 且 $\xi\in\mathbb R^{2n}$。（注意，对
$x,\xi\in\mathbb R^{2n}$，有 $dJ(x)\xi\in\End(\mathbb R^{2n})$。）因为 $J,u,v$
都连续可微，函数 $C:B_\varepsilon\to\End(\mathbb R^{2n})$ 连续，函数
$z\mapsto J(u(z))$ 也连续可微。

因此已经证明 $w:B_\varepsilon\to\mathbb R^{2n}$ 满足定理 2.3.5 的假设，其中把 $J(z)$
替换为 $J(u(z))$。所以每个 $z_0\in B_\varepsilon$ 都有一个邻域 $B_\delta(z_0)$，
使得在该邻域中 $w$ 具有形式 $w(z)=\Phi(z)\sigma(z)$；这里
$\sigma:B_\delta(z_0)\to\mathbb C^n$ 全纯，而
$\Phi:B_\delta(z_0)\to\Hom_{\mathbb R}(\mathbb C^n,\mathbb R^{2n})$ 属于 $W^{1,p}$，
因而连续。

令 $Z$ 为所有使 $w$ 在 $z_0$ 无穷阶消失的点 $z_0\in B_\varepsilon$ 组成的集合。它是
开的：若 $z_0\in Z$，且 $\delta,\sigma,\Phi$ 如上，则 $\sigma$ 在 $z_0$ 无穷阶消失；
由于 $\sigma$ 全纯，这蕴含 $\sigma$、从而 $w$ 在 $B_\delta(z_0)$ 上消失，于是
$B_\delta(z_0)\subset Z$。它也是闭的：若 $z_\nu\in Z$ 收敛到
$z_0\in B_\varepsilon$，并取如上的 $\delta,\sigma,\Phi$，则 $\sigma$ 在一个收敛到
$z_0$ 的点列上消失；所以 $\sigma$、从而 $w$ 在 $z_0$ 无穷阶消失，即 $z_0\in Z$。
因此 $Z$ 是 $B_\varepsilon$ 的既开又闭子集。由于 $0\in Z$，可知
$Z=B_\varepsilon$，从而 $w$ 在 $B_\varepsilon$ 上消失。
\end{proof}

\begin{proof}[定理 2.3.5 的证明]
证明分三步。

\medskip\noindent
\textbf{步骤 1：只需考虑
$(\mathbb R^{2n},J(z))=(\mathbb C^n,i)$ 的情形。}

由假设，存在函数
$\Psi\in W^{1,p}(B_\varepsilon,\Hom_{\mathbb R}(\mathbb C^n,\mathbb R^{2n}))$，
使 $\Psi(z)$ 可逆并把 $J(z)$ 共轭为 $i$，即对每个 $z\in B_\varepsilon$ 都有
$\Psi(z)^{-1}J(z)\Psi(z)=i$。由 $u=\Psi v$ 定义的相应函数 $v$ 满足
\begin{align*}
  0
  &=\partial_su+J\partial_tu+Cu\\
  &=(\partial_s\Psi)v+\Psi\partial_sv+J(\partial_t\Psi)v
    +J\Psi(\partial_tv)+C\Psi v\\
  &:=\Psi\bigl(\partial_sv+i\partial_tv+\widetilde Cv\bigr),
\end{align*}
其中
\begin{equation}\label{eq:2-3-5}
  \widetilde C:=
  \Psi^{-1}(\partial_s\Psi+J\partial_t\Psi+C\Psi)
  \in L^p(B_\varepsilon,\End_{\mathbb R}(\mathbb C^n)).
\end{equation}
（这里 $\End_{\mathbb R}(\mathbb C^n)$ 表示 $\mathbb C^n$ 的实线性变换空间。）因此，
假定 $J=i$ 时结论成立，就可\BookRecordIndex{T0326}推出 $v=\widetilde\Phi(z)\sigma$，其中
$\widetilde\Phi$ 是 $i$-线性的，而 $\sigma$ 全纯。于是
$u=\Phi\sigma$，其中 $\Phi:=\Psi\widetilde\Phi$ 把 $J$ 共轭为 $i$。

\medskip\noindent
\textbf{步骤 2：还可以假定 $C\in\End_{\mathbb R}(\mathbb C^n)$ 是 $i$-线性的。}

可以用任何满足
\[
  A(z)v(z)=\widetilde C(z)v(z)
\]
的 $L^p$ 函数 $A$ 代替 \eqref{eq:2-3-5} 中的函数 $\widetilde C$。为找到合适的 $A$，
把 $\widetilde C$ 分解成复线性部分和复反线性部分：
\[
  \widetilde C=\widetilde C^++\widetilde C^-,
  \qquad
  \widetilde C^\pm:=\frac12(\widetilde C\mp i\widetilde C i).
\]
于是函数
\[
  A:=\widetilde C^++\widetilde C^-D
\]
复线性，当且仅当 $D$ 复反线性。因此要寻找一个不必连续的映射
$D\in L^\infty(B_\varepsilon,\End_{\mathbb R}(\mathbb C^n))$，使 $D(z)$ 复反线性且对
每个 $z$ 都有 $D(z)v(z)=v(z)$。下式给出这样一个映射：
\[
  D(z)\zeta:=
  \begin{cases}
    |v(z)|^{-2}v(z)v(z)^T\overline\zeta,&v(z)\ne0,\\
    0,&v(z)=0.
  \end{cases}
\]

\medskip\noindent
\textbf{步骤 3：当 $J(z)=i$ 且 $C$ 为 $i$-线性时，结论成立。}

必须找出一族复线性自同构 $\Phi(z)$，使函数
$\sigma(z):=\Phi^{-1}(z)u(z)$ 全纯。为此，考虑单位圆盘 $B_1$ 上的\BookRecordIndex{T0210}柯西--黎曼算子。记
\[
  V:=\mathbb C^{n\times n},
  \qquad W:=\mathbb R^{n\times n},
\]
并注意 $W$ 是 $V$ 的拉格朗日子空间。对 $\delta\in(0,\varepsilon)$，考虑函数
$C_\delta\in L^p(B_1,V)$：
\[
  C_\delta(z):=
  \begin{cases}
    C(z),&z\in B_\delta,\\
    0,&z\in B_1\setminus B_\delta.
  \end{cases}
\]
记
\[
  W^{1,p}(B_1,V;W):=
  \{\Phi\in W^{1,p}(B_1,V)\mid\Phi(\partial B_1)\subset W\},
\]
并考虑\BookRecordIndex{T0210}柯西--黎曼算子
\[
  D_\delta:W^{1,p}(B_1,V;W)\longrightarrow L^p(B_1,V),
\]
其定义为
\[
  D_\delta\Phi:=\partial_s\Phi+i\partial_t\Phi+C_\delta\Phi.
\]
算子
\[
  W^{1,p}(B_1,V;W)\longrightarrow L^p(B_1,V),
  \qquad \Phi\longmapsto\partial_s\Phi+i\partial_t\Phi
\]
是满射，其核由常值函数 $\Phi:B_1\to W$ 组成（见定理 C.4.1）。因此算子
\[
  W^{1,p}(B_1,V;W)\longrightarrow L^p(B_1,V)\times W,
  \quad
  \Phi\longmapsto(\partial_s\Phi+i\partial_t\Phi,\Phi(1))
\]
是双射。由于
\[
  \lim_{\delta\to0}\|C_\delta\|_{L^p}=0,
\]
当 $\delta$ 充分小时，算子 $\Phi\mapsto(D_\delta\Phi,\Phi(1))$ 仍是满射。因此对充分小的
$\delta$，存在唯一的 $\Phi_\delta\in W^{1,p}(B_1,V;W)$，满足
\[
  D_\delta\Phi_\delta=0,
  \qquad \Phi_\delta(1)=\mathbf1.
\]
由唯一性，当 $\delta\to0$ 时，$\Phi_\delta$ 收敛到常值函数 $\mathbf1$。所以对充分小的
$\delta$，$\Phi_\delta$ 的值都在 $V$ 的可逆元素中。现在用方程
$u=: \Phi_\delta\sigma$ 定义 $\sigma:B_\varepsilon\to\mathbb C^n$。与步骤 1 相同可知，
$\sigma$ 在 $B_\delta$ 中全纯。
\end{proof}

\begin{exercise}\label{ex:2-3-6}
利用 $2$-球面上的\BookRecordIndex{T0210}柯西--黎曼算子给出步骤 3 的证明。这样的证明不需要边界\BookRecordIndex{T0329}椭圆正则性。
\end{exercise}
